\section{Methods}
\label{sec:methods}

\subsection{Three-valued outcomes}
Every authority decision was three-valued rather than binary: a claim could be promoted, blocked, or left \emph{undetermined}. An undetermined outcome was recorded separately from an outcome determined to be false, and neither category was merged with the other. This distinction makes a missing prerequisite observable without treating the claim as refuted.

\subsection{Prospective design and protected split}
The study protocol, comparator mechanisms, ablations, evaluator, statistical rules, and split-generation procedure were fixed before outcome-bearing evaluation. A public-seed preflight exercised 42 hostile authority tests and the complete 420-case procedure with ten comparator mechanisms, eight ablations, five deterministic repeats, protected scoring, typed-panel validation, and a separately implemented checker. A subsequent isolated evaluation generated a fresh secret-seeded split; the seed was unavailable to every scored mechanism.

The final split contains 420 mechanical-gold cases across 13 families: 60 clean positives and 30 cases in each of 12 hostile or insufficient-evidence families. A 150-case post-evaluation public slice contains 30 clean and 120 hostile cases, while distinct hostile and holdout slices remained protected during scoring. Table~\ref{tab:p4_t1} summarizes the design.

\begin{table}[t]
\centering\scriptsize
\caption{Protected safety battery and custody. Hidden labels remained evaluator-only until all mechanisms were scored.}\label{tab:p4_t1}
\begin{tabular}{lrrrr}\toprule
Slice & Cases & Clean & Hostile & Candidate labels \\\midrule
Post-run public slice & 150 & 30 & 120 & hidden during execution \\
Protected hostile & 120 & 0 & 120 & hidden \\
Protected holdout & 150 & 30 & 120 & hidden \\\midrule
Total & 420 & 60 & 360 & hidden \\\bottomrule
\end{tabular}\end{table}


\subsection{Gold-blind evaluation and custody}
Each scored mechanism received only an opaque case identifier and the candidate-visible evidence packet. Before evaluation, an automated check confirmed that attack family, protected gold, custody class, and expected outcome were absent. Candidate and comparator mechanisms ran in isolated environments without access to study source materials. Five deterministic repeats were monitored for protected-identifier disclosure and external network connections. The evaluator accessed gold only after all scored mechanisms had finished, and a separately implemented checker recomputed the headline counts. These controls establish gold isolation and implementation separation within one project; they do not constitute institutionally independent evaluation or external replication.

\subsection{Comparator panel, ablations, and statistics}
Seven primary comparator mechanisms mirror published ideas: ProvenanceGuard-style source routing, AttributionBench-style multi-source attribution, FIRE-style iterative retrieve/verify, CLAIM-BENCH-style claim evidence, ProvenAI-style citation fidelity/influence, RewardHackingAgents-style contamination/tamper checks, and DeepSciVerify-style evidence escalation. Three auxiliary controls test citation presence, pooled semantic support, and auditability/provenance.

These are protocol-matched reimplementations under one candidate-visible packet and common accounting, not executions of the cited authors' external software.

The eight registered ablations remove, one at a time, exact content binding, source/provenance distinction, checker-lineage independence, hostile-checker discrimination, behavioral influence, evaluator protection/telemetry, or the search-contamination block, or replace the fail-closed rule with a six-of-nine soft-confidence threshold.

The primary endpoint is false authority-promotion rate over 360 gold outcomes that do not permit promotion. The primary hypothesis requires an absolute improvement beyond 0.05 against the strongest frozen comparator. Clean coverage is the promotion rate over 60 clean positives and is required to be non-inferior within 0.05. A third registered contrast tests correct selection of the undetermined outcome. Rate intervals use Wilson 95\% intervals; paired effects use a fixed-seed percentile bootstrap. Every false promotion, false block, abstention, candidate-caused failure, and clean false negative remains in its eligible denominator.
